\chapter{Background and Related Work}
\label{chap:background}

\section{Animal movement data and the missing path}

\begin{itemize}
  \item Animal-borne sensors record discrete, noisy observations of a continuous movement process.
  \item Missing or sparse fixes arise through sampling design, battery limits, transmission loss, environmental obstruction, device orientation, or failed positioning.
  \item Straight-line interpolation supplies one geometric connection but suppresses uncertainty and may cross impossible terrain.
  \item A useful workflow must distinguish:
  \begin{itemize}
    \item observation error from uncertainty about the unobserved path;
    \item short irregularities from genuinely unsupported long gaps;
    \item finite-interval space use from long-term home range;
    \item plausible model realizations from claims about the animal's actual route.
  \end{itemize}
  \item Use movement-ecology literature to connect the movement process, internal state, navigation capacity, and external environment \parencite{nathan2008movement}.
\end{itemize}

\section{Environmental annotation and habitat selection}
\label{sec:environmental-annotation-and-habitat-selection}

\begin{itemize}
  \item Environmental covariates can describe the terrain at observed fixes, candidate endpoints, rasterized steps, or every intermediate model state.
  \item These choices answer different questions and must not be collapsed into one generic notion of ``environment-aware movement''.
  \item Step-selection functions compare used steps with available alternatives and provide the conceptual bridge for estimating relative terrain preference \parencite{thurfjell2014stepselection,avgar2016issa}.
  \item Integrated and state-dependent step-selection approaches motivate separating baseline movement availability from habitat-dependent weighting \parencite{pohle2024states}.
  \item The thesis adapts that used-versus-available logic to terrain-pair crossings inside local movement kernels; the concrete estimator is developed in \cref{sec:ssf-terrain-weights}.
\end{itemize}

\section{Segmentation and behavioral structure}
\label{sec:segmentation-and-behavioral-structure}

\begin{itemize}
  \item Segmentation partitions a trajectory into internally coherent phases; behavioral annotation assigns a label or latent state; gap splitting prevents unsupported steps from contaminating both.
  \item Criteria-based segmentation expresses domain knowledge through temporal and geometric predicates and stable criteria \parencite{alewijnse2014segmentation}.
  \item Gurarie et al. compare exploratory methods including BCPA, Bayesian partitioning, first-passage-time analysis, and multi-state movement models \parencite{gurarie2016behaviour}.
  \item Key implication from Gurarie et al. for this thesis:
  \begin{itemize}
    \item methods encode different assumptions and can identify different kinds or scales of behavioral change;
    \item autocorrelation can create spurious change points when ignored;
    \item the number and interpretation of behaviors depend on the scientific question;
    \item no single automatic annotation should be treated as biological ground truth.
  \end{itemize}
  \item HMM state numbers are statistical labels; ecological names such as resting, foraging, or traveling require post-fit interpretation from feature distributions and context \parencite{zucchini2016hmm}.
\end{itemize}

\section{From Brownian bridges to conditioned discrete walks}

\begin{itemize}
  \item Brownian and dynamic Brownian bridge models provide important reference approaches for uncertainty between observed endpoints \parencite{horne2007bbmm,kranstauber2012dbbmm}.
  \item Their strengths include a clear stochastic interpretation, established usage, and direct bridge-based utilization distributions.
  \item The discretized conditioned-random-walk framework supports a broader family of transition rules:
  \begin{itemize}
    \item Brownian, biased, correlated, mixed, state-dependent, and terrain-dependent motion;
    \item exact endpoint conditioning on a finite grid;
    \item dynamic-programming computation of densities followed by backward sampling;
    \item utilization distributions and visit probabilities under the same model.
  \end{itemize}
  \item The 2024 paper supplies the original framework \parencite{buchin2024randomwalks}; Hagedoorn's thesis supplies the broader algorithmic and case-study context \parencite{hagedoornPhD}; the 2026 manuscript supplies environment-aware extensions \parencite{buchin2026advanced}.
  \item This thesis uses that theory as one component inside a larger application workflow and does not assume the reader has studied those sources.
\end{itemize}

\section{MoveApps and the research gap}

\begin{itemize}
  \item MoveApps makes movement-analysis components composable and accessible, but composition creates strict interface requirements.
  \item A complete workflow must carry coordinate systems, units, timestamps, animal identifiers, states, raster provenance, kernel intervals, and warnings across app boundaries.
  \item Existing theory or individual packages do not by themselves demonstrate that all stages remain compatible on realistic data.
  \item The gap addressed here is therefore practical and scientific:
  \begin{itemize}
    \item connect raw observations and environmental products to the theoretical model;
    \item expose defensible settings instead of implementation-only parameters;
    \item return interpretable artifacts and diagnostics;
    \item evaluate failures and interactions across the complete chain.
  \end{itemize}
\end{itemize}
